% Einheit 1 von 5 – Prozente als Anteile (bank/prozentrechnung/e1.jsonl, zusammenbau v0.1)
%% TODO Zweigzeile Teil 1: was hier gelernt wird (Satz in Schülersprache aus dem Einheitstitel) – Einheit 1
\einheitenkopf[e1]{Einheit 1 von 5 · Prozente als Anteile}
\zweigzeile{ab Kl. 6 (am Gymnasium ab Kl. 5) · P10 oft}
\setcounter{aufgabe}{8}

%% TODO Titel: Ich-kann-Satz fehlt in der Bank (Platzhalter: Kettenname) – Nr. 9
%% TODO Anweisung: ein Satz über der ersten Teilaufgabe fehlt in der Bank – Nr. 9
\begin{aufgabe}{Umwandeln}
\swa{Wie viel Prozent sind grau? \leerfeld[\%]}{\streifen[2]{50}{$0\,\%$}{$100\,\%$}}
\swa{Wie viel Prozent sind grau? \leerfeld[\%]}{\streifen[4]{25}{$0\,\%$}{$100\,\%$}}
\swa{Wie viel Prozent sind grau? \leerfeld[\%]}{\streifen[5]{60}{$0\,\%$}{$100\,\%$}}
\swa{Wie viel Prozent sind grau? \leerfeld[\%]}{\streifen{30}{$0\,\%$}{$100\,\%$}}
\swa{Wie viel Prozent sind grau? \leerfeld[\%]}{\streifen{90}{$0\,\%$}{$100\,\%$}}
\swfrage{Was bleibt gleich, was ändert sich?}
%% TODO Darstellung für schwach fehlt in der Bank (prozentrechnung-e1-k1-s2-v1; Abschnitt „Für schwache Schüler“)
\swz{$\frac{37}{100}$ – wie viel Prozent?}{}
%% TODO Darstellung für schwach fehlt in der Bank (prozentrechnung-e1-k1-s3-v1; Abschnitt „Für schwache Schüler“)
\swz{$\frac{11}{20}$ – wie viel Prozent?}{}
%% TODO Darstellung für schwach fehlt in der Bank (prozentrechnung-e1-k1-s4-v1; Abschnitt „Für schwache Schüler“)
\swz{$0{,}64$ – wie viel Prozent?}{}
%% TODO Darstellung für schwach fehlt in der Bank (prozentrechnung-e1-k1-s5-v1; Abschnitt „Für schwache Schüler“)
\swz{Jeder hundertste Besucher bekommt einen Preis – wie viel Prozent?}{}
%% TODO Darstellung für schwach fehlt in der Bank (prozentrechnung-e1-k1-s6-v1; Abschnitt „Für schwache Schüler“)
\swz{$35\,\%$ der Kinder fahren Rad, die anderen gehen zu Fuß – wie viel Prozent gehen zu Fuß?}{}
\begin{teile}
\teil $5\,\%$ der Pakete kommen zu spät – welche Aussage passt? Kreuze an. \hfill (P10 2020 OS)\\ \kreuz{$5$ von $10$ Paketen kommen zu spät.}\\ \kreuz{Jedes 5. Paket kommt zu spät.}\\ \kreuz{$5$ von $100$ Paketen kommen zu spät.}
\end{teile}
\swa{Schulweg der Kinder einer Schule. Genau eine Aussage ist falsch. Kreuze sie an und schreibe sie richtig. \hfill (P10 2023 OS)\\ \kreuz{Aussage 1: Ein Viertel der Kinder fährt mit dem Bus.}\\ \kreuz{Aussage 2: Die Hälfte der Kinder fährt mit dem Rad.}}{\sachtabelle{lr}{Schulweg & Kinder}{Bus & 150\\ Rad & 200\\ zu Fuß & 180\\ Auto & 70\\ gesamt & 600}}
\swa{Strom einer Gemeinde nach Energiequelle – trage den fehlenden Anteil ein. \hfill (P10 2021 OS)}{\sachtabelle{lr}{Energiequelle & Anteil}{Wind & 38\,\%\\ Sonne & 24\,\%\\ Biogas & \leerzelle\\ Wasser & 9\,\%\\ Sonstiges & 3\,\%\\ gesamt & 100\,\%}}
\end{aufgabe}

%% TODO Titel: Ich-kann-Satz fehlt in der Bank (Platzhalter: Kettenname) – Nr. 10
%% TODO Anweisung: ein Satz über der ersten Teilaufgabe fehlt in der Bank – Nr. 10
\begin{aufgabe}{Prozentangaben ordnen}
%% TODO Darstellung für schwach fehlt in der Bank (prozentrechnung-e1-k2-s1-v1; Abschnitt „Für schwache Schüler“)
\swz{$0{,}4$; $\frac{1}{4}$; $35\,\%$; $\frac{3}{10}$ – der Größe nach, kleinste zuerst?}{}
\end{aufgabe}

%% TODO Titel: Ich-kann-Satz fehlt in der Bank (Platzhalter: Kettenname) – Nr. 11
%% TODO Anweisung: ein Satz über der ersten Teilaufgabe fehlt in der Bank – Nr. 11
\begin{aufgabe}{Umwandeln – Fehler finden, Begründen, Darstellungswechsel, Anwendung}
\swz[3]{Tom: „Jede fünfundzwanzigste Schraube ist fehlerhaft, das sind $25\,\%$.“ – Fehler? Wie heißt es richtig?}{}
\swfrage{Warum ist $\frac{1}{4}$ dasselbe wie $25\,\%$?}
\swa{Färbe $30\,\%$ des Streifens.}{\streifenleer}
\swz[3]{Abstimmung über den Wandertag: In der 7a sind $18$ von $25$ Kindern für den Zoo, in der 7b $21$ von $30$. In welcher Klasse ist der Anteil für den Zoo größer?}{}
\end{aufgabe}

\uebersichtskasten{Prozent heißt Hundertstel: 1 \% = 1/100 = 0,01. \\ Bruch → Prozent: auf Nenner 100 erweitern oder Dezimalzahl mit Komma zwei Stellen nach rechts. \\ 3/4 = 75/100 = 75 \% \quad 0,08 = 8 \% \quad „jeder Fünfte“ = 1/5 = 20 \% \\ Alle Anteile zusammen sind 100 \%.}
